% -*- coding: utf-8 -*-
% 基于所提供的南开大学实验报告模板；编译器：XeLaTeX。
% 保留提供模板的版心、封面字体、校徽、水印和页眉页脚；正文由Markdown生成。
\documentclass[UTF8,a4paper,10pt]{ctexart}
\usepackage[left=3.17cm,right=3.17cm,top=2.74cm,bottom=2.74cm]{geometry}
\usepackage{amsmath,amssymb,graphicx,float,booktabs,longtable,array,tabularx}
\usepackage{caption,multirow,xcolor,listings,enumitem,needspace}
\usepackage{fancyhdr,tikz,eso-pic}
\usepackage[colorlinks,linkcolor=black,urlcolor=blue,citecolor=black]{hyperref}
% 有 Noto 字体时使用模板字体，否则使用 ctex 默认字体。
\IfFontExistsTF{Noto Serif CJK SC}{\setCJKmainfont{Noto Serif CJK SC}}{}
\newcommand{\yihao}{\fontsize{26pt}{36pt}\selectfont}
\newcommand{\erhao}{\fontsize{22pt}{28pt}\selectfont}
\newcommand{\xiaoer}{\fontsize{18pt}{18pt}\selectfont}
\newcommand{\sanhao}{\fontsize{16pt}{24pt}\selectfont}
\newcommand{\xiaosan}{\fontsize{15pt}{22pt}\selectfont}
\newcommand{\sihao}{\fontsize{14pt}{21pt}\selectfont}
\newcommand{\banxiaosi}{\fontsize{13pt}{19.5pt}\selectfont}
\newcommand{\xiaosi}{\fontsize{12pt}{18pt}\selectfont}
\newcommand{\dawuhao}{\fontsize{11pt}{11pt}\selectfont}
\newcommand{\wuhao}{\fontsize{10.5pt}{15.75pt}\selectfont}
\newcommand{\tabincell}[2]{\begin{tabular}{@{}#1@{}}#2\end{tabular}}
\ctexset{
 section={name={,、},number=\chinese{section}},
 subsection={name={（,）},number=\chinese{subsection}},
 subsubsection={name={,.},number=\arabic{subsubsection}}
}
\pagestyle{fancy}
\fancyhf{}
\fancyhead[L]{\kaishu\small Lab2：物理内存与页表}
\fancyhead[R]{\kaishu\small 操作系统实验报告}
\fancyfoot[C]{\thepage}
\setlength{\headheight}{15pt}
\renewcommand{\headrulewidth}{0.1pt}
\renewcommand{\footrulewidth}{0pt}
\newcommand{\HRule}{\rule{\linewidth}{0.5mm}}
\newcommand{\HRulegrossa}{\rule{\linewidth}{1.2mm}}
\definecolor{nkupurple}{RGB}{112,35,84}
\definecolor{codebg}{RGB}{247,248,250}
\lstset{
 escapeinside={(*@}{@*)},
 breaklines=true,basicstyle=\ttfamily\footnotesize,
 keywordstyle=\color{blue!70}\bfseries,
 commentstyle=\color{gray},stringstyle=\ttfamily,
 columns=fullflexible,keepspaces=true,showstringspaces=false,
 numbers=left,numberstyle=\tiny\color{blue!50},numbersep=6pt,
 frame=trbl,rulecolor=\color{gray!40},backgroundcolor=\color{codebg},
 xleftmargin=1.5em,tabsize=4,aboveskip=0.8em,belowskip=0.8em
}
\setlength{\parindent}{2em}
\setlength{\parskip}{0.25em}
\setlength{\emergencystretch}{3em}
\setlist{nosep,leftmargin=2em}
% 透明紫色水印：12% 不透明度，置于前景以透过代码背景显示。
\newcommand{\watermark}{\AddToShipoutPictureFG{%
 \AtPageLowerLeft{\includegraphics[width=\paperwidth,height=\paperheight]{NKU-watermark.pdf}}}}
\hypersetup{pdftitle={Lab2：物理内存与页表}}

\usepackage{xurl}
\urlstyle{tt}

\IfFontExistsTF{KaiTi}{\setCJKfamilyfont{zhkai}[AutoFakeBold=2]{KaiTi}}{}
